% SEGMENT OLP-0583-S01
\documentclass[../../../include/open-logic-section]{subfiles}

\begin{document}

\olfileid{sth}{cardinals}{zfc}
\olsection{$\ZFC$: ஒரு மைல்கல்}

நன்கு வரிசைப்படுத்தல் அடிகோளையும் சேர்த்ததால், கோட்பாட்டின் இறுதி
மைல்கல்லை அடைந்துள்ளோம். $\ZFC$ க்குத் தேவையான எல்லா அடிகோள்களும்
இப்போது நம்மிடம் உள்ளன. விரிவாகக் கூறினால்:

\begin{defn}
$\ZFC$ கோட்பாடு பின்வரும் அடிகோள்களைக் கொண்டது: விரிவுறு
சமத்துவம், ஒன்றிப்பு, இணைகள், அடுக்குக் கணங்கள், முடிவிலி,
அடித்தளம், நன்கு வரிசைப்படுத்தல், மேலும் பிரிப்பு மற்றும்
மாற்றீட்டுத் திட்டவடிவங்களின் எல்லா நிகழ்வுகளும்.
வேறு விதமாகச் சொன்னால், $\ZFC$ என்பது $\ZF$ உடன்
நன்கு வரிசைப்படுத்தலைச் சேர்க்கிறது.
\end{defn}

$\ZFC$ என்பது \emph{செர்மேலோ--ஃபிராங்கெல்} கணக் கோட்பாட்டுடன்
\emph{தேர்வையும்} சேர்த்ததைக் குறிக்கிறது. நாம் சேர்த்த அடிகோளின்
பெயர் ``தேர்வு'' அல்ல, ``நன்கு வரிசைப்படுத்தல்'' என்பதால் இது சற்று
விசித்திரமாகத் தோன்றலாம். ஆனால் பின்னர் {தேர்வை} வகுக்கும்போது,
$\ZF$ இன் பின்னணியில் நன்கு வரிசைப்படுத்தலும் தேர்வும் சமானமானவை
என்பது தெரியவரும் (\olref[choice][woproblem]{thmwochoice} ஐப் பார்க்கவும்).
ஆகவே, இவற்றில் எதை ``அடிப்படை'' அடிகோளாக எடுத்துக்கொள்வது என்பது
முடிவை மாற்றாது. ``$\ZFC$'' என்ற பெயர் இலக்கியத்தில் முற்றிலும்
வழக்கமானது.

\end{document}
